\clearpage
\subsection{How Sage implements the finite predicates}
\label{subsec:rotated-pentagon-sage-source}

The preceding proof has already reduced the theorem to exact finite
predicates.  This subsection makes the correspondence between those
predicates and the Sage source directly inspectable.  The excerpts are copied
in cleaned form from the executable named above: diagnostics, return-value
bookkeeping, and repeated type conversions are omitted, short local names are
used, and long expressions are line-broken.  Named utility constructors and
diagnostics are not repeated.  The source identified by the digest in the
retained transcript remains the runnable owner.

\subsubsection*{Facet duals and the fixed-word quadratic program}
The first excerpt constructs the ten dual vertices in the coordinate order
\((q_1,q_2,p_1,p_2)\), implements \(\omega_0\), and forms the KKT system.
The last four rows impose
\(\sum_i\beta_i a_i=0\), while the final row imposes
\(\sum_i\beta_i=1\).  The function \texttt{q\_value} is exactly the ordered
quadratic form \(Q_\sigma\) used above.

{\small
\begin{verbatim}
def cos_theta():
    return lift(1 - t**2) / lift(1 + t**2)

def sin_theta():
    return lift(2 * t) / lift(1 + t**2)

def rotate(point):
    x, y = point
    c, s = cos_theta(), sin_theta()
    return (c * x - s * y, s * x + c * y)

def pentagon_normals():
    return [(lift(cos_units(5 + 4*k)),
             lift(sin_units(5 + 4*k))) for k in range(5)]

def dual_vertices():
    height = lift(cos_units(2))  # cos(pi/5)
    normals = pentagon_normals()
    vertices = []
    for x, y in normals:
        vertices.append((x / height, y / height, 0, 0))
    for normal in normals:
        x, y = rotate(normal)
        vertices.append((0, 0, x / height, y / height))
    return vertices

def omega(u, v):
    return (u[0] * v[2] - u[2] * v[0]
            + u[1] * v[3] - u[3] * v[1])

def kkt_matrix(sigma):
    m = len(sigma)
    mat = Matrix(F, m + 5, m + 5, 0)
    rhs = vector(F, [0] * (m + 5))
    for i in range(m):
        for j in range(i + 1, m):
            value = omega(DUALS[sigma[i]], DUALS[sigma[j]])
            mat[i, j] = value
            mat[j, i] = value
    for i in range(m):
        for d in range(4):
            value = DUALS[sigma[i]][d]
            mat[i, m + d] = value
            mat[m + d, i] = value
        mat[i, m + 4] = 1
        mat[m + 4, i] = 1
    rhs[m + 4] = 1
    return mat, rhs

def q_value(sigma, beta):
    q = 0
    for i in range(1, len(sigma)):
        for j in range(i):
            q += (beta[i] * beta[j]
                  * omega(DUALS[sigma[j]], DUALS[sigma[i]]))
    return reduced(q)
\end{verbatim}
}

\subsubsection*{The finite block family}
The next excerpt is the implementation of
\(\bigl([Q\mid QQ][P\mid PP]\bigr)^k\), \(k=2,3\).  A block is a singleton
or an ordered adjacent pair.  The selections are disjoint.  The first selected
\(q\)-block fixes the cyclic start, while the remaining relative orders are
all enumerated.

{\small
\begin{verbatim}
def adjacent_same_factor(a, b):
    return (a - b) % 5 in {1, 4}

def blocks(facets):
    result = [(i,) for i in facets]
    for i, j in combinations(facets, 2):
        if adjacent_same_factor(i % 5, j % 5):
            result.append((i, j))
            result.append((j, i))
    return result

def non_overlapping(selection):
    used = set()
    for block in selection:
        if used & set(block):
            return False
        used.update(block)
    return True

def non_overlapping_selections(block_list, k):
    for selection in combinations(block_list, k):
        if non_overlapping(selection):
            yield selection

def enumerate_k_bounce_sigmas(k):
    q_blocks = blocks(Q_FACETS)
    p_blocks = blocks(P_FACETS)
    for qs in non_overlapping_selections(q_blocks, k):
        for ps in non_overlapping_selections(p_blocks, k):
            for q_order in permutations(range(k - 1)):
                for p_order in permutations(range(k)):
                    sigma = list(qs[0]) + list(ps[p_order[0]])
                    for r in range(1, k):
                        sigma.extend(qs[1 + q_order[r - 1]])
                        sigma.extend(ps[p_order[r]])
                    yield tuple(sigma)
\end{verbatim}
}

\subsubsection*{Exact sign cells and transition pruning}
For a rational function in \(K(t)\), Sage isolates every numerator and
denominator root in the open half-domain.  The resulting cells have constant
sign, so one exact algebraic midpoint decides each cell.  The endpoint label
retained in each certificate is diagnostic only; endpoints are handled by the
continuity argument above.  Transition pruning uses the open-cell status and
checks that every mixed transition has no interior sign cut.

{\small
\begin{verbatim}
def real_roots_in_open_half_domain(poly):
    roots = []
    for root, multiplicity in poly.change_ring(AA).roots():
        root = AA(root)
        if AA(0) < root < HALF_DOMAIN_ENDPOINT:
            roots.append(root)
    return sorted(set(roots))

def sign_certificate(expr):
    expr = reduced(expr)
    numerator = expr.numerator()
    denominator = expr.denominator()
    if numerator == 0:
        return SignCertificate("zero", 0, "zero")
    cuts = sorted(set(
        real_roots_in_open_half_domain(numerator)
        + real_roots_in_open_half_domain(denominator)))
    points = [AA(0)] + cuts + [HALF_DOMAIN_ENDPOINT]
    num_aa = aa_polynomial(numerator)
    den_aa = aa_polynomial(denominator)
    signs = []
    for left, right in zip(points, points[1:]):
        sample = (left + right) / 2
        signs.append((num_aa(sample) / den_aa(sample)).sign())
    endpoint = endpoint_label(num_aa, den_aa)
    if all(sign > 0 for sign in signs):
        return SignCertificate("positive_open", len(cuts), endpoint)
    if all(sign < 0 for sign in signs):
        return SignCertificate("negative_open", len(cuts), endpoint)
    return SignCertificate("mixed_open", len(cuts), endpoint)

def open_domain_cells(expressions):
    cuts = []
    for expr in expressions:
        expr = reduced(expr)
        cuts += real_roots_in_open_half_domain(expr.numerator())
        cuts += real_roots_in_open_half_domain(expr.denominator())
    points = [AA(0)] + sorted(set(cuts))
    points += [HALF_DOMAIN_ENDPOINT]
    return [(a + b) / 2 for a, b in zip(points, points[1:])]

def facet_intersection_nonempty(i, j):
    if i == j:
        return True
    if (i < 5) != (j < 5):
        return True
    return adjacent_same_factor(i % 5, j % 5)

def transition_table_open():
    table = {}
    for i in range(10):
        for j in range(10):
            cert = sign_certificate(omega(DUALS[i], DUALS[j]))
            if (i < 5) == (j < 5):
                assert cert.status == "zero"
                table[i, j] = facet_intersection_nonempty(i, j)
            else:
                assert cert.status in {
                    "positive_open", "negative_open"
                }
                assert cert.interior_cut_count == 0
                table[i, j] = (facet_intersection_nonempty(i, j)
                               and cert.status == "positive_open")
    return table

def transition_pruned_sigmas_open():
    table = transition_table_open()
    sigmas = []
    for k in (2, 3):
        for sigma in enumerate_k_bounce_sigmas(k):
            edges = zip(sigma, sigma[1:] + sigma[:1])
            if all(table[i, j] for i, j in edges):
                sigmas.append(sigma)
    return sigmas
\end{verbatim}
}

\subsubsection*{Fail-closed branch classification}
The classifier solves the displayed KKT equations over \(K(t)\).  It accepts
an identically zero action gap as a tie; otherwise, on every cell where all
weights and \(Q_\sigma\) are positive, it requires a strictly positive action
gap.  A singular or nonpositive-gap case not covered by a proved exit is sent
to manual review.  The full loop rejects that status.
For readability, \texttt{status(...)} below abbreviates the executable's
result constructor.

{\small
\begin{verbatim}
ACCEPTED_STATUSES = {
    "no_kkt_solution", "singular_kkt_forced_zero_beta",
    "zero_q_identity", "zero_gap_identity",
    "not_feasible_on_open_domain",
    "strict_gap_positive_on_feasible_open_domain",
}

def forced_zero_weight(beta, kernel):
    for i in range(len(beta)):
        if beta[i] == 0 and all(reduced(v[i]) == 0 for v in kernel):
            return True
    return False

def classify_sigma(sigma, min_action):
    mat, rhs = kkt_matrix(sigma)
    try:
        solution = mat.solve_right(rhs)
    except ValueError:
        augmented = mat.augment(rhs.column())
        if augmented.rank() > mat.rank():
            return status("no_kkt_solution")
        return status("requires_manual_review")

    beta = [reduced(solution[i]) for i in range(len(sigma))]
    q = q_value(sigma, beta)
    kernel = None
    if q == 0:
        kernel = mat.right_kernel().basis()
        if kernel:
            if forced_zero_weight(beta, kernel):
                return status("singular_kkt_forced_zero_beta")
            return status("requires_manual_review")
        return status("zero_q_identity")

    gap = reduced(1 / (2 * q) - min_action)
    if gap == 0:
        return status("zero_gap_identity")

    feasible_cells = 0
    for sample in open_domain_cells(tuple(beta) + (q, gap)):
        feasible = all(sign_at(b, sample) > 0 for b in beta)
        feasible = feasible and sign_at(q, sample) > 0
        if feasible:
            feasible_cells += 1
            if sign_at(gap, sample) <= 0:
                return status("requires_manual_review")

    if feasible_cells == 0:
        if kernel is None:
            kernel = mat.right_kernel().basis()
        if kernel:
            if forced_zero_weight(beta, kernel):
                return status("singular_kkt_forced_zero_beta")
            return status("requires_manual_review")
        return status("not_feasible_on_open_domain")
    return status("strict_gap_positive_on_feasible_open_domain")

for sigma in transition_pruned_sigmas_open():
    classification = classify_sigma(sigma, minimum_action())
    assert classification.status in ACCEPTED_STATUSES
\end{verbatim}
}

In the runnable source every unresolved singular case becomes manual review.
Hence the last assertion passes only when every word reaches one of the six
exits in Table~\ref{tab:rotated-pentagon-certificate-outcomes}.

This comparison of formulas and source is the theorem-specific implementation
audit.  At the software-foundation level, the calculation relies on SageMath
to perform its standard exact number-field arithmetic, linear algebra, and
real-root isolation correctly.  SageMath~\cite{SageMath107} is routinely used
for such exact computer-algebra calculations; verifying SageMath itself is
outside the scope of this thesis.
